\subsection{Методы получения первого приближения}

Все три метода разбираются на одном примере.

\begin{example}\label{ex:transport}
Тарифы $c_{ij}$, запасы $a_i$ и потребности $b_j$:
\begin{center}
\begin{tabular}{c|cccc|c}
\toprule
      & $B_1$ & $B_2$ & $B_3$ & $B_4$ & $a_i$ \\
\midrule
$A_1$ & 14 & 28 & 21 & 28 & 27 \\
$A_2$ & 10 & 17 & 15 & 24 & 20 \\
$A_3$ & 14 & 30 & 25 & 21 & 43 \\
\midrule
$b_j$ & 33 & 13 & 27 & 17 & 90 \\
\bottomrule
\end{tabular}
\end{center}
Задача закрытая: $\sum a_i = 27 + 20 + 43 = 90 = 33 + 13 + 27 + 17 = \sum b_j$.
\end{example}

\subsubsection*{1) Метод северо-западного угла}

Заполняем левую верхнюю («северо-западную») свободную клетку на максимально
возможный объём $\min\{a_i, b_j\}$, после чего вычёркиваем исчерпавшуюся строку
или столбец и повторяем.

\medskip
\noindent Шаги:
\[
  \begin{array}{llll}
    (1,1): & \min\{a_1, b_1\} = 27, &\qquad (2,3): & \min\{a_2, b_3\} = 1, \\
    (2,1): & \min\{a_2, b_1\} = 6,  &\qquad (3,3): & \min\{a_3, b_3\} = 26, \\
    (2,2): & \min\{a_2, b_2\} = 13, &\qquad (3,4): & \min\{a_3, b_4\} = 17 .
  \end{array}
\]

\begin{center}
\begin{tabular}{c|cccc|l}
\toprule
      & $B_1$ & $B_2$ & $B_3$ & $B_4$ & остаток $a_i$ \\
\midrule
$A_1$ & \tc{14}{27} & \tcx{28} & \tcx{21} & \tcx{28} & $27,\ 0$ \\
$A_2$ & \tc{10}{6}  & \tc{17}{13} & \tc{15}{1} & \tcx{24} & $20,\ 14,\ 1,\ 0$ \\
$A_3$ & \tcx{14}    & \tcx{30} & \tc{25}{26} & \tc{21}{17} & $43,\ 17,\ 0$ \\
\midrule
остаток $b_j$ & $33,6,0$ & $13,0$ & $27,26,0$ & $17,0$ & \\
\bottomrule
\end{tabular}
\end{center}

Стоимость полученного плана:
\[
  14 \cdot 27 + 10 \cdot 6 + 17 \cdot 13 + 15 \cdot 1 + 25 \cdot 26 + 21 \cdot 17
  = 1681 .
\]

\subsubsection*{2) Метод минимального элемента}

Заполняем ту свободную клетку, в которой тариф $c_{ij}$ минимален, на объём
$\min\{a_i, b_j\}$; дальше так же.

\medskip
\noindent Шаги:
\[
  \begin{array}{llll}
    (2,1): & \min\{20, 33\} = 20, &\qquad (3,4): & \min\{17, 43\} = 17, \\
    (1,1): & \min\{13, 27\} = 13, &\qquad (3,3): & \min\{26, 13\} = 13, \\
    (1,3): & \min\{14, 27\} = 14, &\qquad (3,2): & \min\{13, 13\} = 13 .
  \end{array}
\]

\begin{center}
\begin{tabular}{c|cccc|l}
\toprule
      & $B_1$ & $B_2$ & $B_3$ & $B_4$ & остаток $a_i$ \\
\midrule
$A_1$ & \tc{14}{13} & \tcx{28} & \tc{21}{14} & \tcx{28} & $27,\ 14,\ 0$ \\
$A_2$ & \tc{10}{20} & \tcx{17} & \tcx{15} & \tcx{24} & $20,\ 0$ \\
$A_3$ & \tcx{14}    & \tc{30}{13} & \tc{25}{13} & \tc{21}{17} & $43,\ 26,\ 13,\ 0$ \\
\midrule
остаток $b_j$ & $33,13,0$ & $13,0$ & $27,13,0$ & $17,0$ & \\
\bottomrule
\end{tabular}
\end{center}

Стоимость:
\[
  14 \cdot 13 + 21 \cdot 14 + 10 \cdot 20 + 30 \cdot 13 + 25 \cdot 13 + 21 \cdot 17
  = 1748 .
\]

\begin{idea}
Здесь метод минимального элемента дал план \emph{дороже}, чем северо-западный угол
($1748$ против $1681$). Это не ошибка: оба метода строят лишь \emph{начальное}
приближение, и гарантий оптимальности ни у одного нет. На шаге~2 тариф $14$
встречается дважды — в клетках $(1,1)$ и $(3,1)$; выбрана первая, и дальше
дешёвые клетки первой строки оказались «съедены», а остаток пришлось везти по
тарифам $30$ и $25$. Обычно минимальный элемент всё же ближе к оптимуму, но
решающий шаг — не выбор начального плана, а последующие итерации метода
потенциалов.
\end{idea}

\subsubsection*{3) Метод двойного предпочтения}

На лекции не разбирался (см. методички).

\begin{idea}
Кратко: в каждой строке отмечают клетку с минимальным тарифом, затем то же самое
делают в каждом столбце. Клетки, отмеченные \emph{дважды} (минимум и по строке,
и по столбцу), заполняют первыми, потом отмеченные один раз, потом остальные —
в каждом случае на объём $\min\{a_i, b_j\}$. Метод — компромисс между дешевизной
северо-западного угла и качеством минимального элемента.
\end{idea}

Далее рассмотрим метод поиска оптимального решения — \textbf{метод (разности)
потенциалов}.
